\documentclass[10pt,a4paper]{amsart}
\usepackage[margin=19mm]{geometry}
\usepackage{amsmath,amssymb,mathtools}
\usepackage[scr=rsfs,cal=euler]{mathalfa}
\usepackage{xfrac}
\usepackage[unicode,hidelinks,bookmarks=false]{hyperref}
\newcommand{\ie}{i.e.\ }
\newcommand{\cf}{cf.\ }
\newcommand{\resp}{respectively\ }
\newcommand{\Subsection}{\subsection}
\providecommand{\nobreakdash}{\nobreak\hbox{-}}
\setlength{\emergencystretch}{2em}
\multlinegap0pt
\usepackage{amssymb}
\usepackage{mathtools}
\usepackage{tikz}
\usepackage{tcolorbox}
\newtheorem{theorem}{Theorem}
\newtheorem{definition}{Definition}
\title{Monadic and cylindric expansions of bounded implication algebras}
\author{Joseph McDonald}
\date{Source: arXiv:2606.03400v1 (CC BY 4.0)}
\begin{document}
\maketitle

\noindent\emph{Source and licence.} Monadic and cylindric expansions of bounded implication algebras. Version arXiv:2606.03400v1 (CC BY 4.0), distributed by arXiv under the Creative Commons Attribution 4.0 International licence, http://creativecommons.org/licenses/by/4.0/, as recorded in the arXiv licence metadata for this exact version and shown on the abstract page https://arxiv.org/abs/2606.03400v1. Full licence notice: \texttt{licenses/arxiv-2606.03400-CC-BY-4.0.txt}.\par\smallskip

\noindent\textit{Editorial scope.} Counted unit: the article's theorem mia (source0.tex lines 301--303). The article's complete original proof of it is delivered with it (source0.tex lines 304--326, one \texttt{proof} environment of 23 lines). No mathematics is written, repaired or re-derived: the statement and the proof are the article's own lines, in the article's own notation, and no retained line is substituted or re-flowed.  Retained uncounted as the source-local context the proof needs: lines 127--129; lines 190--200; lines 235--241.  Nothing was omitted from any retained span, so the package carries no omission record.  No retained line contains a citation, so the wrapper carries no bibliography and no bibliography entry.  No retained line contains a figure environment, an image-inclusion command or drawing code, so no image is delivered and the package declares no asset dependency.  Editorial formatting only, in addition: an AMS \texttt{amsart} preamble that restates the article's own declarations and macros used by the retained lines, namely the environments \texttt{theorem}, \texttt{definition} on separate counters, together with the standard packages those lines need.  The retained environments print in this document as Theorem 1, Definition 1 in order of appearance, whereas the article prints the same items with the numbering of its own full document.  The complete native source of the article as distributed in the arXiv e-print for this exact version is bundled unchanged as \texttt{sources/201969/source0.tex}.
\begin{theorem}\label{ba to bia}
    If $B$ is a Boolean algebra, $\langle B;\pi_c\rangle$ is a bounded implication algebra under $\pi_c(x,y):=x'\vee y$. 
\end{theorem}

\begin{definition}\label{monadic ortholattice}
A \emph{monadic Boolean algebra} is an algebra $\langle A;\wedge,\vee,',0,1,\exists\rangle$ such that: 
\begin{enumerate}
    \item $\langle A;\wedge,\vee,',0,1\rangle$ is a Boolean algebra;
    \item $\exists\colon A\to A$ is a unary operator, known as a \emph{quantifier}, satisfying:
        \begin{enumerate}
        \item $\exists 0=0$ and $x\leq \exists x$;
        \item $\exists(x\wedge\exists y)=\exists x\wedge\exists y$. 
    \end{enumerate}
\end{enumerate}
\end{definition}

\begin{definition}\label{mbia}
    A \emph{monadic implication algebra} is a bounded implication algebra $A$ equipped with an additional operator $\nabla\colon A\to A$ satisfying the following conditions: 
    \begin{enumerate} 
        \item $\nabla c_0=c_0$ and $x\cdot\nabla x=c_1$;
        \item $\nabla((x\cdot(\nabla y\cdot c_0))\cdot c_0)=((\nabla x\cdot(\nabla y\cdot c_0))\cdot c_0)$. 
    \end{enumerate}
\end{definition}

\begin{theorem}\label{mba to mia}
    Every monadic Boolean algebra can be converted into a monadic implication algebra. 
\end{theorem}

\begin{proof}
    By Theorem \ref{ba to bia}, every Boolean algebra $B$ gives rise to a bounded implication algebra $\langle B;\pi_c,0\rangle$ and hence it suffices to demonstrate that the quantifier $\exists\colon B\to B$ satisfies conditions 1 and 2 of Definition \ref{mbia}. For condition 1, observe that $\exists 0=0$ is trivial and that since $\exists$ is increasing and $\vee$ is monotone, we have $\pi_c(x,\exists x)=x'\vee\exists x=1$. For condition 2, if suffices to demonstrate the following: 
    \begin{equation}
\exists\pi_c(\pi_c(x,\pi_c(y,0)),0)=\pi_c(\pi_c(\exists x,\pi_c(\exists y,0)),0).
    \end{equation}
    The first part of calculation proceeds by a repeated application of the definition of $\pi_c$ along with the fact that $\pi_c(x,0)=x'$ for all $x\in B$:
    \begin{equation}
        \exists(\pi_c(\pi_c(x,\pi_c(\exists y,0)),0)=\exists(\pi_c(\pi_c(x,(\exists y)')),0)=\exists(\pi_c(x,(\exists y)')')=\exists(x'\vee(\exists y)')'
    \end{equation}
    De Morgan's identities, the fact that $'$ is an involution, Definition \ref{monadic ortholattice}(2(b)) gives:
    \begin{equation}
        \exists(x'\vee(\exists y)')'=\exists(x''\wedge(\exists y)'')=\exists(x\wedge\exists y)=\exists x\wedge \exists y
    \end{equation}
Applying De Morgan's identities and the fact that $0$ is a unit then yields: 
\begin{equation}
     \exists x\wedge \exists y=((\exists x)'\vee (\exists y)')'=((\exists x)'\vee((\exists y)'\vee 0))'=((\exists x)'\vee((\exists y)'\vee 0))'\vee 0
\end{equation}
Applying the definition of $\pi_c$ then yields the following: 
\begin{equation}
    ((\exists x)'\vee((\exists y)'\vee 0))'\vee 0=\pi_c(\pi_c(\exists x,\pi_c(\exists y,0)),0)
\end{equation}
Equations (5)-(8) then give Equation (4), as desired.    
\end{proof}

\end{document}
